% GENERATED FILE. Edit canonical lessons, then run npm run generate:books.
% canonical-source-hash: fnv1a64:e636b1187e7684ad
% canonical-lessons: UR-C440-jazba, UR-C440-ehsas, UR-C440-mahsus-karna, UR-C440-kaifiyat, UR-C440-mizaj

\chapter{Naming a Feeling}
\label{ch:ur-naming-a-feeling}
\hlchaptermodality{440}

\begin{chapteropening}
\textbf{By the end of this chapter:} \emph{I can name a feeling and say that I felt it.}

Complete the last lesson of chapter 440: I can name a feeling and say that I felt it.
\end{chapteropening}

\section[\texorpdfstring{\ur{جذبہ} (jazbā)}{جذبہ (jazbā)}]{\ur{جذبہ} (jazbā) — an emotion; a strong feeling}

\label{lesson:UR-C440-jazba}



\begin{quote}
Before the new one: say the Urdu for a fruit seller, then the Urdu for a greengrocer.
\end{quote}

\subsection*{You'll want to know: \ur{جذبہ}}
\textbf{\ur{جذبہ}} — \emph{jazbā} — \textquotedblleft{}an emotion; a strong feeling\textquotedblright{}.

\begin{quote}
\ur{ماں} \ur{کا} \ur{پیار} \ur{ایک} \ur{گہرا} \ur{جذبہ} \ur{ہے۔}
\end{quote}

\begin{quote}
\emph{māṅ kā pyār ek gahrā jazbā hai.}
\end{quote}

\begin{quote}
A mother's love is a deep emotion.
\end{quote}

Masculine.

\begin{cousinweb}
\textbf{\ur{جذبہ}} (\emph{jazbā}) is usually traced, through Persian, to Arabic \emph{jadhba}, \textquotedblleft{}a pull, an attraction\textquotedblright{}; the sense \textquotedblleft{}a strong feeling\textquotedblright{} is the one Urdu gives it.
\end{cousinweb}

\begin{grammarlens}[title={What you've built}]
One more word for describing what you lived and felt.
\end{grammarlens}

\subsection*{Your turn}
\begin{itemize}
\raggedright
  \item \emph{Say it:} \emph{jazbā}
  \item \emph{Say it:} \emph{jazbā}, once more
  \item \emph{Say it:} \emph{māṅ kā pyār ek gahrā jazbā hai.} — A mother's love is a deep emotion.
  \item \emph{Recall:} say the Urdu for a fruit seller, then the Urdu for a greengrocer, then say \emph{jazbā} again
\end{itemize}

\begin{tcolorbox}[breakable,colback=teal!4,colframe=teal!35!black,title={Before you move on}]
What does \ur{جذبہ} mean? (\textquotedblleft{}An emotion; a strong feeling\textquotedblright{}.) Say it once more.
\end{tcolorbox}
\section[\texorpdfstring{\ur{احساس} (ehsās)}{احساس (ehsās)}]{\ur{احساس} (ehsās) — a feeling; an awareness}

\label{lesson:UR-C440-ehsas}



\begin{quote}
Before the new one: say the Urdu for a greengrocer, then the Urdu for an emotion.
\end{quote}

\subsection*{You'll want to know: \ur{احساس}}
\textbf{\ur{احساس}} — \emph{ehsās} — \textquotedblleft{}a feeling; an awareness\textquotedblright{}.

\begin{quote}
\ur{مجھے} \ur{سردی} \ur{کا} \ur{احساس} \ur{ہوا۔}
\end{quote}

\begin{quote}
\emph{mujhe sardī kā ehsās huā.}
\end{quote}

\begin{quote}
I became aware of the cold.
\end{quote}

Masculine.

\begin{cousinweb}
\textbf{\ur{احساس}} (\emph{ehsās}) is usually traced, through Persian, to Arabic \emph{iḥsās}, \textquotedblleft{}perceiving, sensing\textquotedblright{}.
\end{cousinweb}

\begin{grammarlens}[title={What you've built}]
One more word for describing what you lived and felt.
\end{grammarlens}

\subsection*{Your turn}
\begin{itemize}
\raggedright
  \item \emph{Say it:} \emph{ehsās}
  \item \emph{Say it:} \emph{ehsās}, once more
  \item \emph{Say it:} \emph{mujhe sardī kā ehsās huā.} — I became aware of the cold.
  \item \emph{Recall:} say the Urdu for a greengrocer, then the Urdu for an emotion, then say \emph{ehsās} again
\end{itemize}

\begin{tcolorbox}[breakable,colback=teal!4,colframe=teal!35!black,title={Before you move on}]
What does \ur{احساس} mean? (\textquotedblleft{}A feeling; an awareness\textquotedblright{}.) Say it once more.
\end{tcolorbox}
\section[\texorpdfstring{\ur{محسوس} \ur{کرنا} (mahsūs karnā)}{محسوس کرنا (mahsūs karnā)}]{\ur{محسوس} \ur{کرنا} (mahsūs karnā) — to sense; to feel (inside)}

\label{lesson:UR-C440-mahsus-karna}



\begin{quote}
Before the new one: say the Urdu for an emotion, then the Urdu for a feeling.
\end{quote}

\subsection*{You'll want to know: \ur{محسوس} \ur{کرنا}}
\textbf{\ur{محسوس} \ur{کرنا}} — \emph{mahsūs karnā} — \textquotedblleft{}to sense; to feel (inside)\textquotedblright{}.

\begin{quote}
\ur{میں} \ur{نے} \ur{اس} \ur{کا} \ur{درد} \ur{محسوس} \ur{کیا۔}
\end{quote}

\begin{quote}
\emph{maiṅ ne us kā dard mahsūs kiyā.}
\end{quote}

\begin{quote}
I felt his pain.
\end{quote}

Past, with \ur{نے} (he or she …): \ur{اس} \ur{نے} \ur{تھکن} \ur{محسوس} \ur{کی}, \emph{us ne thakan mahsūs kī}.

\begin{cousinweb}
\textbf{\ur{محسوس} \ur{کرنا}} (\emph{mahsūs karnā}) is built from \textbf{\ur{محسوس}}, usually traced, through Persian, to Arabic \emph{maḥsūs}, \textquotedblleft{}perceived by the senses\textquotedblright{}, and \textbf{\ur{کرنا}}, \textquotedblleft{}to do\textquotedblright{}.
\end{cousinweb}

\begin{grammarlens}[title={What you've built}]
One more word for describing what you lived and felt.
\end{grammarlens}

\subsection*{Your turn}
\begin{itemize}
\raggedright
  \item \emph{Say it:} \emph{mahsūs karnā}
  \item \emph{Say it:} \emph{mahsūs karnā}, once more
  \item \emph{Say it:} \emph{maiṅ ne us kā dard mahsūs kiyā.} — I felt his pain.
  \item \emph{Recall:} say the Urdu for an emotion, then the Urdu for a feeling, then say \emph{mahsūs karnā} again
\end{itemize}

\begin{tcolorbox}[breakable,colback=teal!4,colframe=teal!35!black,title={Before you move on}]
What does \ur{محسوس} \ur{کرنا} mean? (\textquotedblleft{}To sense; to feel (inside)\textquotedblright{}.) Say it once more.
\end{tcolorbox}
\section[\texorpdfstring{\ur{کیفیت} (kaifiyat)}{کیفیت (kaifiyat)}]{\ur{کیفیت} (kaifiyat) — a state of mind; a state of body}

\label{lesson:UR-C440-kaifiyat}



\begin{quote}
Before the new one: say the Urdu for a feeling, then the Urdu for to sense.
\end{quote}

\subsection*{You'll want to know: \ur{کیفیت}}
\textbf{\ur{کیفیت}} — \emph{kaifiyat} — \textquotedblleft{}a state of mind; a state of body\textquotedblright{}.

\begin{quote}
\ur{آج} \ur{اس} \ur{کی} \ur{کیفیت} \ur{بہتر} \ur{ہے۔}
\end{quote}

\begin{quote}
\emph{āj us kī kaifiyat behtar hai.}
\end{quote}

\begin{quote}
She is in a better state today.
\end{quote}

Feminine.

\begin{cousinweb}
\textbf{\ur{کیفیت}} (\emph{kaifiyat}) is usually traced, through Persian, to Arabic \emph{kayfiyya}, \textquotedblleft{}how a thing is\textquotedblright{}; Urdu uses it for the state a person is in.
\end{cousinweb}

\begin{grammarlens}[title={What you've built}]
One more word for describing what you lived and felt.
\end{grammarlens}

\subsection*{Your turn}
\begin{itemize}
\raggedright
  \item \emph{Say it:} \emph{kaifiyat}
  \item \emph{Say it:} \emph{kaifiyat}, once more
  \item \emph{Say it:} \emph{āj us kī kaifiyat behtar hai.} — She is in a better state today.
  \item \emph{Recall:} say the Urdu for a feeling, then the Urdu for to sense, then say \emph{kaifiyat} again
\end{itemize}

\begin{tcolorbox}[breakable,colback=teal!4,colframe=teal!35!black,title={Before you move on}]
What does \ur{کیفیت} mean? (\textquotedblleft{}A state of mind; a state of body\textquotedblright{}.) Say it once more.
\end{tcolorbox}
\section[\texorpdfstring{\ur{مزاج} (mizāj)}{مزاج (mizāj)}]{\ur{مزاج} (mizāj) — a mood; a temperament}

\label{lesson:UR-C440-mizaj}



\begin{quote}
Before the new one: say the Urdu for to sense, then the Urdu for a state of mind.
\end{quote}

\subsection*{You'll want to know: \ur{مزاج}}
\textbf{\ur{مزاج}} — \emph{mizāj} — \textquotedblleft{}a mood; a temperament\textquotedblright{}.

\begin{quote}
\ur{آج} \ur{ابو} \ur{کا} \ur{مزاج} \ur{اچھا} \ur{ہے۔}
\end{quote}

\begin{quote}
\emph{āj abbū kā mizāj achchhā hai.}
\end{quote}

\begin{quote}
Dad is in a good mood today.
\end{quote}

Masculine.

\begin{cousinweb}
\textbf{\ur{مزاج}} (\emph{mizāj}) is usually traced, through Persian, to Arabic \emph{mizāj}, \textquotedblleft{}a mixture\textquotedblright{}, and so the blend that makes up a person's temper.
\end{cousinweb}

\begin{grammarlens}[title={What you've built}]
One more word for describing what you lived and felt.
\end{grammarlens}

\subsection*{Your turn}
\begin{itemize}
\raggedright
  \item \emph{Say it:} \emph{mizāj}
  \item \emph{Say it:} \emph{mizāj}, once more
  \item \emph{Say it:} \emph{āj abbū kā mizāj achchhā hai.} — Dad is in a good mood today.
  \item \emph{Recall:} say the Urdu for to sense, then the Urdu for a state of mind, then say \emph{mizāj} again
\end{itemize}

\begin{tcolorbox}[breakable,colback=teal!4,colframe=teal!35!black,title={Before you move on}]
What does \ur{مزاج} mean? (\textquotedblleft{}A mood; a temperament\textquotedblright{}.) Say it once more.
\end{tcolorbox}
